\documentclass[11pt,a4paper]{article}
% GitHub-Overleaf sync marker

\usepackage[margin=2.5cm]{geometry}
\usepackage{amsmath}
\usepackage{booktabs}
\usepackage{tabularx}
\usepackage{siunitx}
\usepackage{graphicx}
\usepackage{float}
\usepackage[hidelinks]{hyperref}

\title{Intelligent Multi-Sensor Attitude Determination and Reliability Assessment for a CubeSat}
\author{Antonio Grigaliunas}
\date{}

\begin{document}

\maketitle

% ============================================================
% MA4822 marking structure:
% Motivation: Users and Their Needs .............. 10 marks
% Literature Survey .............................. 10 marks
% Description of Proposed Solution ............... 50 marks
% Implementation ................................ 10 marks
% Results ....................................... 10 marks
% Conclusion and Future Works .................... 5 marks
% References ..................................... 5 marks
% ============================================================

\section{Motivation: Users and Their Needs}

Small satellites require knowledge of their orientation in space in order to perform mission functions such as pointing antennas, orienting payloads, managing power generation, and supporting autonomous attitude-control functions. For a CubeSat, this attitude information must be obtained using sensors that are compact, low-power, low-cost, and sufficiently robust for autonomous operation.

The primary users considered in this project are CubeSat mission designers and operators who require a simple attitude-determination subsystem that can estimate the spacecraft orientation and, at the same time, indicate whether that estimate should be trusted. This second requirement is important because an attitude estimate can become inaccurate when a sensor is noisy, biased, disturbed, or fails completely.

The proposed system therefore addresses two user needs:
\begin{enumerate}
    \item estimate the three-dimensional CubeSat attitude using low-cost complementary sensors; and
    \item provide an interpretable reliability score that warns when the attitude estimate is likely to be degraded.
\end{enumerate}

The project is sensor-centric because the autonomy of the proposed system depends directly on the quality, fusion, and interpretation of the Sun-sensor and magnetometer measurements.

\section{Literature Survey}

\subsection{CubeSat Attitude Determination}

Attitude determination is a central function of a spacecraft attitude determination and control system (ADCS), because pointing and control actions depend on knowledge of the spacecraft orientation. In vector-based attitude determination, onboard measurements of known physical directions are compared with the corresponding directions expressed in a reference frame. Common reference directions include the Sun, the Earth's magnetic field, and stars. For full three-dimensional attitude determination, at least two non-parallel vector observations are required: one vector fixes a direction in space, while the second resolves the remaining rotation about that direction \cite{shuster1981triad}.

This principle is particularly attractive for CubeSats, where mass, volume, power, and cost strongly constrain the available sensors. Nurgizat et al.\ demonstrated a low-cost CubeSat orientation system based solely on solar sensors and a three-axis magnetometer, using solar sensors on the six spacecraft faces to recover the satellite--Sun vector and the magnetometer to obtain the geomagnetic-field vector \cite{nurgizat2023lowcost}. Their work shows that a Sun--magnetic sensor pair can provide a practical alternative to more complex attitude sensors for missions that require coarse-to-moderate attitude knowledge. In-orbit work on the UVSQ-SAT CubeSat has also demonstrated the practical use of TRIAD for spacecraft attitude determination, reporting sunlight attitude accuracy better than approximately \SI{3}{\degree} at one standard deviation \cite{finance2021uvsqsat}. These studies motivate the use of a Sun-sensor and magnetometer combination in the present project.

\subsection{Coarse Sun Sensors}

Sun sensors provide a direct observation of the Sun direction relative to the spacecraft body. Coarse Sun sensors are especially suitable for small satellites because they can be implemented with simple photoelectric devices and require relatively little mass, power, and processing. For an approximately cosine-response sensor, the measured signal varies with the projection of the Sun vector onto the sensor normal. Multiple sensors placed on different faces can therefore be combined to reconstruct the three-dimensional Sun direction.

A six-face arrangement is directly relevant to CubeSats. Nurgizat et al.\ placed solar sensors on the $+X$, $-X$, $+Y$, $-Y$, $+Z$, and $-Z$ faces and processed their illumination measurements to estimate the satellite--Sun vector \cite{nurgizat2023lowcost}. NASA's Small Spacecraft Technology State-of-the-Art report also lists compact cosine Sun sensors for small spacecraft; for example, the Space Micro CSS-01/CSS-02 has a \SI{120}{\degree} full-cone field of view and an accuracy of approximately \SI{5}{\degree} (3$\sigma$) \cite{nasa2024soa}. These characteristics support the use of a coarse cosine-response model in this project rather than a fine Sun sensor model.

The principal limitation is that Sun-sensor accuracy depends on illumination geometry and that direct Sun information is unavailable during eclipse. The UVSQ-SAT flight study illustrates this limitation: attitude estimation based on the available vector measurements degraded substantially in eclipse compared with sunlight operation \cite{finance2021uvsqsat}. The present project therefore focuses on sunlit operation and treats eclipse handling as future work.

\subsection{Magnetometers}

A three-axis magnetometer provides a measurement of the local geomagnetic-field vector in the spacecraft body frame. When the spacecraft position and an Earth magnetic-field model are available, the corresponding reference-frame magnetic direction can be predicted and used as an attitude reference. Magnetometers are attractive for CubeSats because they are small, low-power devices and provide an absolute vector measurement independent of the Sun.

The PNI RM3100 is used as the representative magnetometer in this project. Colagrossi et al.\ investigated an effective sensor architecture for full-attitude determination in the HERMES nanosatellites and reported RM3100 performance values including a worst-axis noise level of \SI{17}{\nano\tesla}, bias of \SI{440}{\nano\tesla}, and scale-factor error of 4885 ppm \cite{colagrossi2023hermes}. These values provide a physically grounded basis for the nominal magnetometer model.

A magnetometer is also vulnerable to magnetic fields generated by the spacecraft itself. Electrical currents, magnetic actuators, and other onboard components can contaminate the ambient geomagnetic measurement. Hoffmann and Moldwin showed that spacecraft-generated magnetic signals can be significant in CubeSat measurements \cite{hoffmann2022spacecraftnoise}. Consequently, the present model separates the small intrinsic magnetometer noise from a larger externally imposed magnetic-disturbance term and evaluates the disturbance independently.

\subsection{Vector-Based Attitude Determination}

The TRIAD algorithm is one of the earliest and simplest methods for determining three-axis attitude from two vector observations. Black introduced the passive two-vector attitude-determination approach in 1964 \cite{black1964triad}, and Shuster and Oh later presented TRIAD together with QUEST in a systematic comparison of attitude determination from vector observations \cite{shuster1981triad}. TRIAD constructs orthonormal bases from two reference vectors and their corresponding measured body-frame vectors, after which the attitude matrix follows directly from the transformation between the two bases.

TRIAD is deterministic and computationally inexpensive, which makes it attractive for a small embedded system. Its main limitation is that it is not an optimal weighted estimator and its result depends more strongly on the vector treated as the primary vector. Shuster and Oh contrasted this with QUEST, which solves a weighted attitude-estimation problem using an arbitrary number of observations \cite{shuster1981triad}. More sophisticated recursive estimators, such as an extended Kalman filter, can additionally combine measurements over time and incorporate spacecraft dynamics. For example, the UVSQ-SAT study compared TRIAD with a multiplicative extended Kalman filter and found that the recursive estimator offered advantages under more difficult measurement conditions \cite{finance2021uvsqsat}. Nevertheless, because the present project uses exactly two absolute vector observations and emphasizes sensing, fault behaviour, and interpretability rather than high-performance dynamic estimation, TRIAD provides an appropriate baseline.

\subsection{Sensor Reliability and Fuzzy Logic}

Attitude estimation alone does not indicate whether the underlying sensor measurements are trustworthy. Sensor faults such as bias, drift, excessive noise, and complete signal loss can propagate into the estimated state and degrade the performance of an autonomous system. Reviews of sensor fault diagnosis therefore emphasize early fault identification and the use of analytical or intelligent methods to distinguish healthy from abnormal measurements \cite{li2020sensorfault}. This is particularly relevant for a CubeSat because simply adding redundant hardware increases mass, power consumption, and system complexity.

Fuzzy logic provides a convenient way of combining several uncertain indicators without requiring a precise probabilistic model of every possible fault. Fuzzy inference systems have previously been applied to sensor-fault detection by converting measurement discrepancies or residuals into fuzzy conditions and combining them through interpretable rules \cite{khan2012fuzzyfault}. Such an approach is well suited to the current problem because the available indicators---for example vector-consistency error, magnetometer magnitude error, Sun-sensor reconstruction residual, and the number of healthy sensors---do not individually provide a definitive measure of attitude quality.

The proposed work therefore extends the conventional Sun--magnetic attitude-determination architecture by adding a separate fuzzy reliability layer. The attitude estimate itself is still obtained using TRIAD, while the fuzzy system evaluates measurement quality and produces a reliability score. In this way, the project does not attempt to replace the attitude estimator with fuzzy logic; instead, fuzzy reasoning is used as an interpretable health-assessment mechanism for the sensor-derived attitude solution.

\section{Description of Proposed Solution}

\subsection{System Architecture}

The proposed attitude-determination system combines:
\begin{itemize}
    \item six coarse Sun sensors mounted on the $\pm X$, $\pm Y$, and $\pm Z$ faces of the CubeSat;
    \item one three-axis magnetometer;
    \item a TRIAD attitude-estimation algorithm; and
    \item a fuzzy-logic reliability layer.
\end{itemize}

The sensor measurements are first converted into two estimated body-frame reference directions: the Sun direction $\hat{\mathbf{s}}_b$ and the geomagnetic-field direction $\hat{\mathbf{m}}_b$. Their corresponding reference-frame directions, $\mathbf{s}_r$ and $\mathbf{m}_r$, are assumed known. TRIAD then estimates the body-from-reference rotation matrix $\hat{\mathbf{R}}$.

In parallel, four measurement-quality indicators are evaluated and passed to a fuzzy inference system. The fuzzy system outputs a reliability score between 0 and 100\%, allowing the attitude estimate and its estimated trustworthiness to be reported together.

\subsection{Six-Face Sun-Sensor Model}

Each Sun sensor has a unit surface normal $\mathbf{n}_i$. The ideal normalized response is modelled as
\begin{equation}
    I_i = \max\left(0,\mathbf{n}_i^\mathsf{T}\mathbf{s}_b\right),
\end{equation}
where $\mathbf{s}_b$ is the true Sun direction expressed in the CubeSat body frame. The $\max(0,\cdot)$ operation represents the fact that a sensor on the opposite side of the spacecraft cannot receive direct sunlight.

Measurement imperfections are represented by
\begin{equation}
    y_i = g_i I_i + b_i + n_i,
\end{equation}
where $g_i$ is the sensor gain, $b_i$ is an offset, and $n_i$ is additive Gaussian measurement noise. Negative readings are clipped to zero.

For the six orthogonal sensor normals, the body-frame Sun direction is reconstructed using
\begin{equation}
    \mathbf{s}_{\mathrm{raw}}
    =
    \mathbf{N}^{\mathsf{T}}\mathbf{y},
\end{equation}
where $\mathbf{N}$ contains the six sensor normals as rows and $\mathbf{y}$ contains the six sensor readings. The estimated Sun direction is then
\begin{equation}
    \hat{\mathbf{s}}_b
    =
    \frac{\mathbf{s}_{\mathrm{raw}}}
    {\|\mathbf{s}_{\mathrm{raw}}\|}.
\end{equation}

For the $\pm X$, $\pm Y$, and $\pm Z$ arrangement, this is equivalent to forming the signed difference between opposing face measurements along each body axis.

\subsection{Magnetometer Model}

The three-axis magnetometer measurement is modelled as
\begin{equation}
    \mathbf{m}_{\mathrm{meas}}
    =
    \mathbf{G}\mathbf{m}_b
    +
    \mathbf{b}
    +
    \mathbf{n}
    +
    \mathbf{d},
\end{equation}
where $\mathbf{m}_b$ is the true body-frame magnetic-field direction, $\mathbf{G}$ contains scale-factor errors, $\mathbf{b}$ represents bias, $\mathbf{n}$ is random sensor noise, and $\mathbf{d}$ is an optional magnetic disturbance used in the fault studies.

The estimated magnetic direction is obtained by normalizing the measured vector:
\begin{equation}
    \hat{\mathbf{m}}_b
    =
    \frac{\mathbf{m}_{\mathrm{meas}}}
    {\|\mathbf{m}_{\mathrm{meas}}\|}.
\end{equation}

The intrinsic magnetometer noise and the spacecraft-generated disturbance are treated separately. This distinction is important because spacecraft-generated magnetic interference can be much larger than the intrinsic noise of the sensor.

\subsection{TRIAD Attitude Estimation}

Two non-parallel vectors are used to construct orthonormal triads in the reference and body frames. In the reference frame,
\begin{align}
    \mathbf{t}_{1r} &= \mathbf{s}_r, \\
    \mathbf{t}_{2r} &=
    \frac{\mathbf{s}_r \times \mathbf{m}_r}
    {\|\mathbf{s}_r \times \mathbf{m}_r\|}, \\
    \mathbf{t}_{3r} &= \mathbf{t}_{1r}\times\mathbf{t}_{2r}.
\end{align}

Similarly, using the measured body-frame vectors,
\begin{align}
    \mathbf{t}_{1b} &= \hat{\mathbf{s}}_b, \\
    \mathbf{t}_{2b} &=
    \frac{\hat{\mathbf{s}}_b \times \hat{\mathbf{m}}_b}
    {\|\hat{\mathbf{s}}_b \times \hat{\mathbf{m}}_b\|}, \\
    \mathbf{t}_{3b} &= \mathbf{t}_{1b}\times\mathbf{t}_{2b}.
\end{align}

The two triad matrices are
\begin{equation}
    \mathbf{T}_r =
    \begin{bmatrix}
        \mathbf{t}_{1r} & \mathbf{t}_{2r} & \mathbf{t}_{3r}
    \end{bmatrix},
    \qquad
    \mathbf{T}_b =
    \begin{bmatrix}
        \mathbf{t}_{1b} & \mathbf{t}_{2b} & \mathbf{t}_{3b}
    \end{bmatrix},
\end{equation}
and the estimated body-from-reference rotation matrix is
\begin{equation}
    \hat{\mathbf{R}}
    =
    \mathbf{T}_b\mathbf{T}_r^\mathsf{T}.
\end{equation}

In simulation, the true rotation matrix $\mathbf{R}$ is known, so the principal attitude error can be calculated from
\begin{equation}
    \mathbf{R}_{e}
    =
    \hat{\mathbf{R}}\mathbf{R}^{\mathsf{T}},
\end{equation}
followed by
\begin{equation}
    \theta_e
    =
    \cos^{-1}
    \left(
        \frac{\mathrm{tr}(\mathbf{R}_e)-1}{2}
    \right).
\end{equation}

This true error is used only for simulation validation and is not an input to the onboard reliability estimator.

\subsection{Fault and Degradation Scenarios}

The simulation investigates both nominal sensor imperfections and more severe faults. The main fault scenarios are:
\begin{itemize}
    \item complete loss of one Sun-sensor reading;
    \item external magnetic disturbance;
    \item simultaneous Sun-sensor failure and magnetic disturbance.
\end{itemize}

The six Sun sensors are placed symmetrically on the CubeSat faces, so averaged single-sensor failure effects are expected to be similar when the spacecraft attitudes are sampled uniformly. The all-six failure comparison is therefore primarily used as a consistency check rather than to identify one uniquely critical face.

\subsection{Fuzzy Reliability Assessment}

The reliability system uses four crisp indicators that can be computed from the sensor measurements themselves.

\subsubsection{Sun-Sensor Health Fraction}

If $N_h$ Sun sensors are considered healthy, the health fraction is
\begin{equation}
    H_s = \frac{N_h}{6}.
\end{equation}
This quantity is fuzzified into \emph{poor}, \emph{degraded}, and \emph{good} states.

\subsubsection{Sun--Magnetic Vector Consistency}

The angular separation between the two reference vectors is
\begin{equation}
    \alpha_r =
    \cos^{-1}\left(
        \mathbf{s}_r^\mathsf{T}\mathbf{m}_r
    \right),
\end{equation}
while the measured separation is
\begin{equation}
    \alpha_m =
    \cos^{-1}\left(
        \hat{\mathbf{s}}_b^\mathsf{T}\hat{\mathbf{m}}_b
    \right).
\end{equation}

The consistency error is
\begin{equation}
    e_c = |\alpha_m-\alpha_r|.
\end{equation}
A small value indicates that the relative geometry between the two sensed vectors remains physically consistent.

\subsubsection{Magnetometer Norm Error}

Because the magnetic vector is normalized in the model, an abnormal magnitude can indicate magnetic disturbance. The indicator is
\begin{equation}
    e_m =
    \left|
        \|\mathbf{m}_{\mathrm{meas}}\|-1
    \right|.
\end{equation}

\subsubsection{Sun-Sensor Reconstruction Residual}

The estimated Sun direction is used to reconstruct the expected sensor readings:
\begin{equation}
    \hat{I}_i =
    g_i\max\left(0,\mathbf{n}_i^\mathsf{T}\hat{\mathbf{s}}_b\right)+b_i.
\end{equation}

The root-mean-square reconstruction residual is
\begin{equation}
    e_s =
    \sqrt{
        \frac{1}{N}
        \sum_{i=1}^{N}
        \left(
            y_i-\hat{I}_i
        \right)^2
    }.
\end{equation}
This indicator helps detect inconsistencies such as a failed Sun sensor that may not be obvious from the reconstructed Sun direction alone.

\subsubsection{Fuzzy Inference and Defuzzification}

Triangular and trapezoidal membership functions are used to convert the four crisp indicators into fuzzy beliefs. The rule base follows the MA4822 fuzzy-logic approach: conjunction is evaluated with the MIN operator, while multiple rules supporting the same output class are combined with the MAX operator.

Five output reliability classes are used:
\begin{equation}
    \{\text{Very Low, Low, Medium, High, Very High}\},
\end{equation}
with representative crisp values
\begin{equation}
    \{10,30,50,70,90\}.
\end{equation}

If the aggregated rule strengths are denoted by $P_k$ and their representative reliability values by $V_k$, the final reliability score is obtained using weighted-sum defuzzification:
\begin{equation}
    R_{\mathrm{rel}}
    =
    \frac{
        \sum_k V_k P_k
    }{
        \sum_k P_k
    }.
\end{equation}

Five output levels are used instead of only Low/Medium/High in order to reduce the strong quantization that appears when only three representative output values are used. Some residual banding remains because the output still uses discrete representative values. Continuous output membership functions with centroid defuzzification are therefore retained as a possible future improvement.

\subsection{Sensor Model and Parameter Selection}

The Sun-sensor model uses a baseline normalized random noise standard deviation of 0.02. NASA's reported \SI{5}{\degree} (3$\sigma$) coarse Sun-sensor accuracy can be related approximately to the cosine response near \SI{45}{\degree} incidence:
\begin{equation}
    \sigma_I
    \approx
    \sin(45^\circ)
    \left(
        \frac{5^\circ}{3}
        \frac{\pi}{180^\circ}
    \right)
    \approx 0.0206.
\end{equation}

The magnetometer parameters are based on a PNI RM3100-class device. Published nanosatellite work reports a worst-axis noise level of \SI{17}{\nano\tesla}, bias of \SI{440}{\nano\tesla}, and scale-factor error of 4885 ppm \cite{colagrossi2023hermes}. The Earth's surface magnetic-field intensity is approximately \SIrange{25}{65}{\micro\tesla} \cite{noaaGeomagFAQ}; a representative value of \SI{50}{\micro\tesla} is used for normalization.

Table~\ref{tab:sensor-parameters} summarizes the parameters and explicitly distinguishes sourced quantities from deliberate sensitivity-study assumptions.

\begin{table}[H]
\centering
\caption{Sensor and environment parameters used in the simulation.}
\label{tab:sensor-parameters}
\footnotesize
\renewcommand{\arraystretch}{1.15}
\begin{tabularx}{\textwidth}{@{}p{2.6cm}p{3.1cm}p{3.6cm}X@{}}
\toprule
\textbf{Subsystem} & \textbf{Parameter} & \textbf{Value} & \textbf{Basis} \\
\midrule
Sun sensor
& Architecture
& 6 cosine sensors, $\pm X,\pm Y,\pm Z$
& Proposed coarse six-face CubeSat arrangement. \\

Sun sensor
& Angular accuracy
& \SI{5}{\degree} (3$\sigma$)
& Literature value for Space Micro CSS-01/CSS-02 \cite{nasa2024soa}. \\

Sun sensor
& Normalized random-noise std.
& 0.02
& Derived from the \SI{5}{\degree} (3$\sigma$) value using a cosine-response linearization near \SI{45}{\degree}. \\

Sun sensor
& Per-face gain factors
& \shortstack[l]{$[1.03,\ 0.98,\ 0.96,$\\$1.02,\ 1.05,\ 0.97]$}
& Deliberate residual calibration-mismatch assumption. \\

Sun sensor
& Per-face offsets
& \shortstack[l]{$[0.010,\ 0.005,\ 0,$\\$0.008,\ 0.012,\ 0.004]$}
& Deliberate residual zero-bias assumption. \\

Magnetometer
& Device basis
& PNI RM3100 class
& Representative nanosatellite magnetometer \cite{colagrossi2023hermes}. \\

Magnetometer
& Noise level
& \SI{17}{\nano\tesla}
& RM3100 worst-axis literature value \cite{colagrossi2023hermes}. \\

Magnetometer
& Bias
& \SI{440}{\nano\tesla}
& RM3100 worst-axis literature value \cite{colagrossi2023hermes}. \\

Magnetometer
& Scale-factor error
& 4885 ppm
& RM3100 worst-axis literature value \cite{colagrossi2023hermes}. \\

Environment
& Reference magnetic-field magnitude
& \SI{50}{\micro\tesla}
& Representative value within NOAA's approximate \SIrange{25}{65}{\micro\tesla} surface range \cite{noaaGeomagFAQ}. \\

Magnetometer
& Normalized noise std.
& $3.4\times10^{-4}$
& Derived as $17/50000$. \\

Magnetometer
& Normalized worst-axis bias
& $8.8\times10^{-3}$
& Derived as $440/50000$. \\

Fault study
& Magnetic-disturbance sweep
& \SIrange{0}{7.5}{\micro\tesla}
& Sensitivity sweep reaching the scale of reported CubeSat onboard magnetic interference \cite{hoffmann2022spacecraftnoise}. \\

Monte Carlo
& Samples per case
& 1000
& Simulation-design choice. \\
\bottomrule
\end{tabularx}
\end{table}

\section{Implementation}

The proposed system is implemented in Python using NumPy for numerical operations and Matplotlib for result visualization. The implementation is separated into sensor models, attitude estimation, fuzzy reliability, Monte Carlo analysis, and fault-analysis scripts.

The main software components are:
\begin{itemize}
    \item \texttt{sun\_sensor.py}: six-face cosine Sun-sensor model;
    \item \texttt{magnetometer.py}: three-axis magnetometer model;
    \item \texttt{attitude\_estimation.py}: Sun-vector reconstruction, TRIAD, and attitude-error calculations;
    \item \texttt{fuzzy\_reliability.py}: membership functions, fuzzy rules, and weighted-sum defuzzification;
    \item \texttt{monte\_carlo.py}: random-attitude noise sensitivity study;
    \item \texttt{fault\_analysis.py}: nominal and faulted scenarios;
    \item \texttt{magnetometer\_sweep.py}: magnetic-disturbance sensitivity study.
\end{itemize}

For each Monte Carlo sample, a uniformly distributed random three-dimensional rotation is generated using a normalized random quaternion. The known reference Sun and magnetic vectors are rotated into the body frame, passed through the sensor models, reconstructed, fused using TRIAD, and compared with the known simulated attitude. The same random seed is used when comparing fault scenarios so that differences are caused by the injected fault rather than by a different random attitude sequence.

\section{Results}

\subsection{Nominal Attitude-Estimation Accuracy}

% TODO: Insert final Monte Carlo summary and attitude-error histogram.
% Current physically based model gives approximately 2.3 deg mean
% attitude error at Sun noise = 0.02 and nominal RM3100 noise.

\subsection{Sensor-Noise Sensitivity}

% TODO: Insert attitude_error_vs_sensor_noise.png and explain that
% Sun-sensor noise dominates the tested nominal RM3100 noise range.

\subsection{Fault Analysis}

% TODO: Insert fault_analysis_summary.png and summarize Healthy,
% Sun sensor failed, Mag disturbance, and Combined fault cases.

\subsection{Fuzzy Reliability Validation}

% TODO: Insert fuzzy_reliability_validation.png and discuss the inverse
% relationship between true attitude error and estimated reliability.

\subsection{Magnetic-Disturbance Sweep}

% TODO: Insert magnetometer_disturbance_sweep.png and
% magnetometer_fault_detection.png, with physical disturbance units in uT.

\section{Conclusion and Future Works}

% TODO: Replace with final conclusion after figures/results are frozen.

The proposed architecture demonstrates that a six-face coarse Sun-sensor array and a three-axis magnetometer can be combined to estimate the full CubeSat attitude using TRIAD, while a fuzzy-logic layer can provide an independent indication of estimate reliability. The simulation framework also allows controlled injection of sensor noise, calibration mismatch, sensor failure, and magnetic disturbance.

Future work should include hardware-in-the-loop or laboratory validation, replacement of assumed Sun-sensor calibration errors with measured calibration data, use of orbit- and location-dependent Sun and geomagnetic reference models, explicit eclipse handling, and possible addition of a gyroscope for dynamic propagation between absolute vector measurements. The fuzzy reliability system could also be improved using continuous output membership functions with centroid defuzzification to reduce the remaining score quantization.

\bibliographystyle{ieeetr}
\bibliography{references}

\end{document}
